%% The Elicit Admin manual — UC-029, "Consult the Administrator's Manual".
%%
%% Build it with docs/manual/build-manual.sh, which renders it through a TeX Live container
%% and stamps it with the application version and build date (UC-029 BR-003). See README.md.
%%
%% Every procedure carries the screen it describes (BR-002). The figures come from
%% docs/manual/capture/capture-screenshots.mjs, which walks a greenfield stack; a caption is
%% part of the instruction, so re-read the captions whenever the figures are regenerated.

\documentclass[11pt, a4paper]{article}

\usepackage[a4paper, margin = 25mm, bottom = 28mm]{geometry}
%% The footer's product name; the brand package is the umbrella's, shared by every Elicit manual.
\newcommand{\elicitProduct}{Elicit Admin}
\usepackage{elicit-brand}
\usepackage{hyperref}
\elicitAccent{elicitVerdigrisLight}

\hypersetup{
  pdftitle = {Elicit Admin — Running the Console},
  pdfauthor = {Elicit Software LLC},
  pdfsubject = {Step-by-step instructions for running the Elicit admin console},
  colorlinks = true,
  linkcolor = elicitInkDark,
  urlcolor = elicitInkDark,
}

\graphicspath{{./}}

%% \adminonly — marks a procedure only the administrator role may carry out (UC-029 BR-006, A5).
%% The margin is widened to hold the word: at the default width LaTeX hyphenates it.
\setlength{\marginparwidth}{22mm}
\setlength{\marginparsep}{2mm}
\newcommand{\adminonly}{%
  \marginpar{\raggedright\footnotesize\textcolor{elicitVerdigris}{\textsf{Administrator}}}%
}

\begin{document}

%% ------------------------------------------------------------------ title page
\elicitTitlePage{Running the Console}{The Elicit Admin manual}%
  {Departments \textperiodcentered\ users \textperiodcentered\ subjects and access codes \textperiodcentered\ message templates \textperiodcentered\ survey definitions \textperiodcentered\ participation \textperiodcentered\ the System pages}%
  {This manual was produced by the same build as the application it ships inside, and describes that version.\\[0.5em] \textcopyright\ 2025--2026 Elicit Software LLC \textperiodcentered\ PolyForm Noncommercial License 1.0.0}

%% ------------------------------------------------------------------ contents
\pagenumbering{roman}
\setcounter{page}{1}
\elicitContents

\clearpage
\pagenumbering{arabic}
\setcounter{page}{1}

\elicitpart{Getting started}{What the console is for, the words it uses, and how to sign in.
Read this part first, whatever your role.}

%% ------------------------------------------------------------------ 1
\section{What the console manages}
\label{sec:overview}

This manual is for the two people who use the Elicit admin console: the \term{administrator},
who sets the console up and keeps it running, and the \term{user}, who works in it day to day.
Both read the same manual. Where a procedure is open only to the administrator it is marked in
the margin, so that a user knows what to ask for rather than looking for a screen that is not
there.

It describes the console of an installed system. Installing Elicit — the database, the identity
provider, and the Survey and Admin services — and provisioning it with a brand, or with the
translations of the languages it offers, is the subject of the \emph{installation manual}. Building a
survey is the subject of the \emph{author's manual}. This manual points at both rather than
repeating them.

\subsection{The vocabulary}

The console names eight things, and this manual uses those names and no others.

\begin{itemize}
  \item A \term{department} is the group a subject belongs to. It carries a code, a sender
        address, and the message template sent to its subjects. Every
        console account is granted one or more departments, and sees only their subjects.
  \item A \term{user} is a console account. A \term{role} is what that account may do:
        \val{elicit_admin} or \val{elicit_user}.
  \item A \term{subject} is a person you have registered — a name, an email address, a
        department. A \term{respondent} is that same person as the survey sees them: their
        progress, their answers, the moment they finished.
  \item An \term{access code} is the credential a respondent enters to reach their survey. The
        console issues one when you register the subject. It is never called a token.
  \item A \term{message template} is the subject line and body of an invitation or a reminder.
  \item A \term{survey definition} is the survey itself as a portable file, exported from the
        Author tool and applied here.
\end{itemize}

\subsection{Finding your way around}

The navigation on the left is the whole console. \ui{Search Subjects}, \ui{Register
Subjects} and \ui{Participation} are the day's work and are open to both roles. \ui{Admin} and
\ui{System} hold the administrator's screens and stay collapsed until opened. The header
carries \ui{Manual}, which opens this document in a new tab, and the language selector when the
deployment offers more than one language (Chapter~\ref{sec:language}); \ui{Logout} is at the
foot of the navigation.

\screenshot{01-navigation}{The console with both navigation sections open. Search Subjects,
Register Subjects and Participation are open to every account; the Admin and System sections
are the administrator's.}

\begin{elicitnote}[Where a procedure lives]
\ui{Admin} holds \ui{Departments}, \ui{Message Templates}, \ui{Users}, \ui{Import Clients},
\ui{Import Respondent},
\ui{Apply Survey Definition} and \ui{Export Survey Definition}. \ui{System} holds the six
read-only screens of Chapter~\ref{sec:system}. \ui{Manual (PDF)} opens this document.
\end{elicitnote}

%% ------------------------------------------------------------------ 2
\section{Signing in}
\label{sec:signin}

The console never asks you for a password. It hands you to the identity provider your
deployment was installed with, and that is where you sign in.

\begin{enumerate}
  \item Open the console's address in a browser. It sends you straight on to the identity
        provider.
  \item Sign in with the username and password you were given.
  \item The identity provider returns you to the console.
\end{enumerate}

\screenshot{02-sign-in}{The identity provider's sign-in page. The console itself never
collects a password, so this page — not the console — is where a forgotten password is reset.}

\subsection{When your account has no department}

A console account that is not assigned to a department has nothing to show. Rather than open an
empty console, Elicit holds it behind a notice.

\screenshot{03-department-required}{The notice an account with no department is held
behind. What it offers depends on your role.}

\begin{itemize}
  \item As an \textbf{administrator} it offers \ui{Add a department}, which takes you to
        Chapter~\ref{sec:departments}. The department you create is assigned to you, and the console opens.
  \item As a \textbf{user} it has no such action. Ask an administrator to assign you to a
        department, then sign in again.
\end{itemize}

Both also offer \ui{Logout}, and \ui{Open the manual}, which opens this document.

This is what a freshly installed deployment looks like to its first administrator: no
department exists yet, and creating one is the first thing to do.

\subsection{When the console does not let you in}

Two other screens can meet you after the identity provider has accepted you:

\begin{description}
  \item[“there is no user named … in the application”] The identity provider knows your
        account, but the console has no user record for it, or the record is not active. An
        administrator adds or reactivates it on \ui{Users} (Chapter~\ref{sec:users}).
  \item[\ui{Access Restricted}] Your account holds neither \val{elicit_user} nor
        \val{elicit_admin}. Ask whoever administers the identity provider for the role.
\end{description}

\elicitpart{Setting up the console}{The groundwork an administrator lays before anyone can be
registered: departments, console accounts and the messages subjects receive. Every chapter in
this part is open to the administrator role only.}

%% ------------------------------------------------------------------ 3
\section{Departments}\adminonly
\label{sec:departments}

A department groups the subjects you register. It decides which surveys those subjects may be
registered onto and which message templates they are sent.

A new installation has no department at all, and until one exists nobody can be registered. The
first department is assigned to whoever creates it.

\begin{elicitnote}[The first department and the first template need each other]
A department must name a message template, and a template must belong to a department. On a new
installation, create the department with a placeholder \ui{Default message ID} such as
\val{1}, write its template (Chapter~\ref{sec:templates}), then edit the department and put the
template's real ID in. Use a number: an ID that is not a number makes every registration in
that department fail with an access-code error.
\end{elicitnote}

\subsection{The list}

\ui{Admin} $>$ \ui{Departments} lists every department with its code, the message templates it
sends, its sender address and its surveys.

\screenshot{04-departments}{The Departments screen. \ui{New Department} creates one; the
pencil in the \ui{Edit} column opens an existing one.}

\subsection{Creating one}

\begin{enumerate}
  \item Open \ui{Admin} $>$ \ui{Departments} and choose \ui{New Department}.
  \item Fill in \ui{Department name} — what the department is called throughout the console.
  \item Optionally fill in \ui{Department code}, a short abbreviation such as \val{CARD} or
        \val{HR}. The code is what identifies this department to another Elicit deployment, so
        moving a respondent between deployments (Chapter~\ref{sec:move}) depends on it.
  \item Fill in \ui{Default message ID}: the ID of the message template sent to this
        department's subjects when they are registered. Templates are written in
        Chapter~\ref{sec:templates}; the ID is the \ui{ID} column of that screen, the one after
        \ui{Edit}.
  \item Fill in \ui{From email}, the department's sender address: every invitation and
        reminder to this department's subjects is sent from it. Use an address the mail relay
        is allowed to send for; ask whoever runs the relay if you are not sure, because a relay
        may refuse an address it does not accept, and a recipient's server may file the mail as
        spam.
  \item Check \ui{Surveys}: the surveys this department's subjects may be registered onto. A new
        department starts with every installed survey; remove the ones it does not run. Each
        survey is shown with its ID, which a CSV upload names it by (Section~\ref{sec:csv}).
  \item Choose \ui{Create department}.
\end{enumerate}

\screenshot{05-edit-department}{Creating a department. \ui{Create department} stays disabled
until every required field is valid, so a greyed-out button means a field above it is still
wrong.}

\begin{elicitnote}[Which template, and how many]
\ui{Default message ID} takes more than one ID, separated by commas. Every template listed is
sent, which is how a deployment that must send both an invitation and a consent notice does it.
\end{elicitnote}

\begin{elicitcaution}
The template the ID names has to exist. Nothing on this screen checks that it does, and a
department pointing at a template that was never created sends nothing at all.
\end{elicitcaution}

To change a department later, open it from the pencil in the \ui{Edit} column. The form is the
same one; the button reads \ui{Update department}.

\begin{elicitcaution}[A department registers onto its surveys only]
Nobody can register a subject in a department onto a survey the department does not list: not on
\ui{Register Subjects}, not by CSV, not through the registration API. A survey applied after the
department was created is in none, so add it here before anyone registers onto it
(Chapter~\ref{sec:install}). A department may list no survey at all.
\end{elicitcaution}

%% ------------------------------------------------------------------ 4
\section{Users and roles}\adminonly
\label{sec:users}

A console account is a \term{user}. What it may do is its \term{role}; which subjects it may see
is its department assignment. Both are set here.

\subsection{The list}

\ui{Admin} $>$ \ui{Users} lists every console account, whether it is active, and the departments
it reaches.

\screenshot{06-users}{The Users screen. The note above the grid says which parts of an
account this console controls, which depends on how the deployment authenticates.}

\subsection{Adding or changing one}

\begin{enumerate}
  \item Open \ui{Admin} $>$ \ui{Users} and choose \ui{Add User}, or the pencil on an existing
        row.
  \item Fill in \ui{Username}. It has to match the account in the identity provider — that is
        what ties the two together.
  \item Fill in \ui{First name} and \ui{Last name}.
  \item Leave \ui{Active} ticked. Clearing it keeps the account and its history but stops it
        signing in.
  \item Choose the departments this account may reach under \ui{Departments}. An account may
        hold several.
  \item Choose the \ui{Role}.
  \item \ui{Save}.
\end{enumerate}

\screenshot{07-edit-user}{Editing a console account. \ui{Departments} decides which subjects
the account sees; \ui{Role} decides what it may do with them.}

\subsection{What each role reaches}

\begin{itemize}
  \item \val{elicit_admin} — the whole console: departments, users, message templates,
        survey definitions, the System screens, and the export action of Chapter~\ref{sec:move}.
  \item \val{elicit_user} — \ui{Search Subjects}, \ui{Register Subjects} and
        \ui{Participation}, for the departments the account holds, and this manual.
\end{itemize}

Choose the account's \emph{highest} role: \val{elicit_admin} implies \val{elicit_user}.

A third role, \val{elicit_importer}, is not a user's and is not offered here. It belongs to
the systems that feed subjects in through the registration API, which are entered on a page of
their own (Section~\ref{sec:import-clients}), so that such a system is not also an account that
can read subjects.

\begin{elicitnote}[Where the role is really decided]
Which of these two places is authoritative depends on the deployment's
\opt{elicit.authorization.mode}. Under \val{OIDC} the roles come from the identity
provider and this screen assigns departments only; under \val{DATABASE} the role set here is
the one that counts. The note above the grid tells you which case you are in.
\end{elicitnote}

\subsection{The two seeded accounts}

The database migrations create \val{admin} and \val{user} so that a new installation can be
signed into at all. Until they are renamed the console carries a warning on every screen
(Chapter~\ref{sec:system}). Rename them here, and — under \val{DATABASE} authorization — rename the matching
accounts in the identity provider to the same values.

\subsection{Import clients}
\label{sec:import-clients}

A system that registers subjects through the registration API is an \term{import client}, not a
user: it has no first or last name and never signs in to the console. \ui{Admin} $>$
\ui{Import Clients} lists these systems, whether each is active, and the departments it may
register into.

\screenshot{31-import-clients}{The Import Clients page. Each row is a system that registers
subjects through the registration API, not a person.}

\begin{enumerate}
  \item Open \ui{Admin} $>$ \ui{Import Clients} and choose \ui{Add Import Client}, or the pencil
        on an existing row.
  \item Fill in \ui{Name} with the name the client's tokens carry. For a Keycloak service
        account that is \val{service-account-} followed by the client id.
  \item Add a \ui{Description} if it helps to say what the system is.
  \item Leave \ui{Active} ticked. Clearing it keeps the record and stops the client registering.
  \item Choose under \ui{Departments} the departments the client may register subjects into.
  \item \ui{Save}.
\end{enumerate}

\screenshot{32-edit-import-client}{Adding an import client. \ui{Name} is the name the
client's tokens carry; \ui{Departments} decides where it may register subjects.}

\ui{Delete}, on the form of an existing import client, removes the record after a confirmation.
The subjects it registered are kept.

\begin{elicitnote}[What the record does and does not do]
The client holds \val{elicit_importer} in the identity provider, and no other Elicit role;
this page says only into which departments it registers. The registration API refuses a client
that has the role and no active record here, a client whose token also carries
\val{elicit_admin} or \val{elicit_user}, and every caller without the role, users
included: people register subjects on \ui{Register Subjects}. A name is a user or an import
client, never both, so each of the two screens refuses a name the other holds, and an import
client record gives no access to the console.
\end{elicitnote}

%% ------------------------------------------------------------------ 5
\section{Message templates}\adminonly
\label{sec:templates}

A message template is what a subject receives: a subject line, a body, and the format the body is
written in. Templates belong to a department, and a department sends the one its \ui{Default
message ID} names.

\subsection{The list}

\ui{Admin} $>$ \ui{Message Templates} lists every template with its department, subject line
and format.

\screenshot{08-message-templates}{The Message Templates screen. The \ui{ID} column is what a
department's \ui{Default message ID} refers to.}

\subsection{Writing one}

\begin{enumerate}
  \item Open \ui{Admin} $>$ \ui{Message Templates} and choose \ui{New Message Template}.
  \item Choose the \ui{Department} this template belongs to.
  \item Choose the \ui{MIME type}: \val{text/plain} for a plain message,
        \val{text/html} to write the body as HTML.
  \item Write the \ui{Subject}.
  \item Write the \ui{Body}, and put \val{<ACCESS_CODE>} where the recipient's access code
        should appear.
  \item \ui{Save} creates it. On an existing template the button reads \ui{Update}.
\end{enumerate}

\screenshot{09-edit-message-template}{A template with the access code placeholder in its
body. \ui{Message preview} below the form shows the body as the recipient will receive it.}

\begin{elicitnote}[The placeholder]
\val{<ACCESS_CODE>} is replaced with that subject's own access code as the message is sent —
it is not something you fill in. A template without it still sends, but the recipient has no way
into their survey. When no template at all carries it, \ui{System} $>$ \ui{Overview} reports
that as outstanding setup.
\end{elicitnote}

Give the recipient the survey's address as well as the code. The code alone is not a link.

\elicitpart{Day-to-day work}{Registering subjects, inviting them, following their progress,
printing their reports and watching how a survey is doing. Open to every console account.}

%% ------------------------------------------------------------------ 6
\section{Registering subjects}
\label{sec:register}

Registering a subject records who they are, issues the \term{access code} that reaches their
survey, and queues the invitation: the message templates their department names
(Chapter~\ref{sec:invite}).

\begin{enumerate}
  \item Open \ui{Register Subjects}.
  \item Choose the \ui{Department}, then the \ui{Survey}. The survey list offers only the
        department's surveys (Chapter~\ref{sec:departments}); either is chosen for you when there
        is only one.
  \item Fill in \ui{First name} and \ui{Last name}. \ui{Middle name}, \ui{Date of birth} and
        \ui{Phone} are optional; a phone number is written \texttt{\small \#\#\#-\#\#\#-\#\#\#\#}.
        A name can be up to 50 characters.
  \item Fill in \ui{Email}, which is required, up to 255 characters. The invitation goes to this
        address.
  \item Fill in \ui{External ID} if this subject has an identifier in another system, such as a
        medical record number. It must be unique within the department, and an External ID the
        department has excluded is refused.
  \item \ui{Save}.
\end{enumerate}

\screenshot{10-register-subject}{The registration form, filled in. The panel on the right
documents the two other ways to register: a CSV upload and the REST API.}

The console confirms with \ui{Subject saved} and clears the form for the next one.

\screenshot{11-register-result}{What a successful registration looks like. The access code
issued at this moment is not shown here — read it from \ui{Search Subjects}.}

\subsection{Changing a subject}

Find the subject on \ui{Search Subjects} and use the pencil in the \ui{Edit} column. The
registration form reopens with their details and the button reads \ui{Update subject}. Their
access code does not change.

\subsection{Registering many at once}
\label{sec:csv}

\ui{Upload CSV} takes a file of subjects, one per row, with these columns in order:

\begin{center}\small
\val{surveyId, departmentId, xid, firstName, lastName, middleName, dob, email, phone}
\end{center}

Every row names its survey in the first column. A file in which any row lacks it is rejected
with a warning naming those lines, and nothing is imported; the survey chosen in the form plays
no part in an upload.

\begin{enumerate}
  \item Write the file. The first column is the survey's ID, shown beside each survey's name in
        the department's \ui{Surveys}. The second is the department's numeric ID: open the
        department on \ui{Departments} and read the number at the end of the page's address
        (\val{edit-department/3}). The third, the External ID, may be left empty
        (\val{2,3,,Ana,...}). First name, last name and email are required; a date of birth is
        \val{yyyy-MM-dd} or \val{MM/dd/yyyy}. A row with no phone may end at the email. Rows
        beginning with \texttt{\small\#} are ignored.
  \item Open \ui{Register Subjects} and choose \ui{Upload CSV}.
  \item Pick the file. Nothing is registered yet: a confirmation says how many subjects the file
        will register on each survey and in each department, and lists each row it will refuse
        with its line number and the reason — a department you are not assigned to, a department
        that does not have the row's survey, or a value that fails its rule.
  \item Choose \ui{Import} to register them, or \ui{Cancel} to register no one. The result lists
        each subject with its access code and whether it was a \ui{New Subject}, an
        \ui{Existing Subject} already registered, or an \ui{Excluded Subject}; the refused rows
        are listed under \ui{Errors} with their line numbers.
\end{enumerate}

\screenshot{35-csv-confirm}{The confirmation an upload asks for: how many subjects it will
register on each survey and in each department, and any row it will refuse. Nobody is
registered until \ui{Import} is chosen.}

\screenshot{34-csv-import}{The result of a CSV upload: each subject with its access code and
outcome, and any row that could not be read, with its line number.}

The \ui{CSV file structure} panel on the registration screen repeats the columns with an example.
\ui{REST API instructions} documents the equivalent endpoints for a system that registers
subjects programmatically; such a system is entered as an import client
(Section~\ref{sec:import-clients}), and the endpoints are not open to a user's account.

%% ------------------------------------------------------------------ 7
\section{Inviting and reminding}
\label{sec:invite}

Registering a subject queues their invitation, and the console sends queued mail every five
minutes, so a newly registered subject is invited without anyone doing anything more. A
reminder, or an invitation sent again, goes from the subject's row on \ui{Search Subjects}, one
subject at a time.

\begin{enumerate}
  \item Open \ui{Search Subjects} and find the subject (Chapter~\ref{sec:search}).
  \item In their row's \ui{Action} column, choose \ui{Send email}.
  \item A \ui{Submit} button appears beside it. Choose it.
\end{enumerate}

\screenshot{12-search-email-action}{Choosing \ui{Send email} on a subject's row. Nothing is
sent until \ui{Submit} is chosen.}

The console reports \ui{Email sent successfully}, or tells you it failed.

What is sent is the template — or templates — that the subject's \emph{department} names in its
\ui{Default message ID}, addressed to the subject's own email address, with
\val{<ACCESS_CODE>} replaced by their access code. It is sent from the department's
\ui{From email}; only when that is blank is the service's own sender address used, the one
\ui{System} $>$ \ui{Email} shows.

A reminder is the same action again: choose \ui{Send email} a second time and the same template
goes out with the same code.

\begin{elicitcaution}
A department whose \ui{Default message ID} names no existing template queues nothing, so its
subjects are never invited. \ui{System} $>$ \ui{Overview} reports a deployment with no template
at all.
\end{elicitcaution}

%% ------------------------------------------------------------------ 8
\section{Finding subjects and watching progress}
\label{sec:search}

\ui{Search Subjects} is both the directory and the progress board.

\subsection{Searching}

\ui{Department(s)} is required — it defaults to every department your account reaches. The other
fields narrow the result: \ui{Access code}, \ui{First name}, \ui{Last name}, \ui{Email},
\ui{Phone}. Leave them empty to list everyone. Choose \ui{Search}.

\screenshot{13-search-filters}{The search filters. \ui{Department(s)} carries a required
marker; the rest are optional and combine.}

\subsection{Reading the results}

\screenshot{14-search-results}{A result. \ui{Status} is the subject's progress through their
survey; \ui{Access code} is the credential they enter to reach it.}

Every column sorts. \ui{Page size} at the foot of the grid sets how many rows are shown at once.
The grid refreshes itself every ten seconds, so a subject working through their survey moves
without the page being reloaded.

\ui{Status} has three values:

\begin{itemize}
  \item \textbf{Not Started} — the access code has been issued but never used.
  \item \textbf{In Progress} — they have signed in to the survey and not yet finished it.
  \item \textbf{Finished} — they have completed it. Only then is a report available.
\end{itemize}

\subsection{Auto-registered respondents}

Where the survey site lets anyone start a survey by typing an access code of their own (the
setting \opt{accessCode.autoRegister}), the people who do so were never registered here. They
are listed with \ui{Auto-registered} where the department would be, and with no name or contact
details: only the access code they chose and their status. Every console account sees them,
whatever its departments, and \ui{Department(s)} offers \ui{Auto-registered} to show only them.
Their row has no \ui{Edit} and no \ui{Send email}, because there is no subject to edit and no
address to write to; reports and the export are offered as for anyone else.

%% ------------------------------------------------------------------ 9
\section{Reports}
\label{sec:reports}

A report is generated from a respondent's answers, and only once they have finished.

\begin{enumerate}
  \item Open \ui{Search Subjects} and find the subject. Their \ui{Status} must read
        \textbf{Finished}.
  \item In their row's \ui{Action} column, choose \ui{Print reports}.
  \item Choose \ui{Submit}.
\end{enumerate}

\screenshot{15-report-download}{\ui{Print reports} is offered only on a respondent whose
status is \textbf{Finished}; the other rows offer the actions their status allows.}

The console assembles the report as a PDF and hands it to the browser, then reports \ui{Reports
generated successfully}. Print or save it from the browser's own viewer.

\begin{elicitnote}[Where reports come from]
The reports belong to the survey, not to the console. Each is a \term{report service} the
survey's author listed in the authoring tool, and they arrive with the survey definition: the
built-in respondent report, or a small service the deployment runs, such as the three reports
of the Family Health History Survey. A survey that lists none prints an empty document.
\ui{System} $>$ \ui{Connections} lists the report services this deployment will call and checks
that each one answers.
\end{elicitnote}

If one report fails, the PDF carries an error block in its place and the others print as
usual. A report that the respondent sees on their own report page but that is missing from the
printout is a fault in that report service, not in the console: tell whoever runs it. How a
report service is built and installed is the subject of the installation manual's custom
reports chapters.

%% ------------------------------------------------------------------ 10
\section{Monitoring a survey's participation}
\label{sec:participation}

\ui{Participation} shows how one survey is doing across your departments, as counts. It never
shows a respondent: for one person, use \ui{Search Subjects}.

\begin{enumerate}
  \item Open \ui{Participation}.
  \item Choose the \ui{Survey}. \ui{Departments} starts at \ui{All departments}, which is every
        department your account reaches and the auto-registered respondents
        (Section~\ref{sec:search}); choose one or more to narrow it. The auto-registered get a
        row of their own in the department table when there are any.
  \item Optionally set \ui{Registered from} and \ui{Registered to}. Both days are included.
\end{enumerate}

\screenshot{33-participation}{The Participation page for one survey: the counts and rates, the
same by department and by week, where people stop, and the marked questions. A count from 1 to 4
reads \textbf{< 5}.}

The page recomputes as soon as a filter changes. It has five parts:

\begin{plaintable}{tab:participation}{The five parts of the Participation page}{@{\hspace{4pt}}>{\raggedright\arraybackslash}p{0.2\linewidth} >{\raggedright\arraybackslash}p{0.72\linewidth}@{\hspace{4pt}}}{\thead{Part} & \thead{What it shows}}
\textbf{Participation} & How many subjects were registered, and how many have not started, are
in progress and have finished: the three states \ui{Search Subjects} shows. \ui{Start rate} is
started over registered, \ui{Completion rate} finished over started, and \ui{Overall rate}
finished over registered. \ui{Stalled} counts the subjects in progress who have not opened a page
or saved an answer for some days (14 unless your site changed it): the ones to remind, from
\ui{Search Subjects}. \\
\textbf{By department} & The same, one row per department. \\
\textbf{By week} & Per week, Monday to Sunday, how many were registered, started and finished in
it. \\
\textbf{Progress} & For subjects still in progress, the step of the last page they were shown, so
you can see where people stop. For those who finished, how long each step took them (median and
quartiles) and how long the whole survey took. A subject who started before the deployment
recorded pages is counted as \ui{Not recorded}. \\
\textbf{Marked questions} & One card for each question or section the survey's author marked for
this report. Among finished respondents, \ui{Presented} counts each time a question was shown
(once per repeat), \ui{Answered} how many of those were answered, \ui{Response rate} the share
answered, and \ui{Reach} the share of finished respondents who were shown it at all. A choice
question lists its options with a count and share each, a number question its minimum, quartiles
and maximum, and a date question its range and a count per month. Free-text answers are counted,
never shown. \\
\end{plaintable}

\begin{elicitnote}[Why some counts read “< 5”]
A count from 1 to 4 is shown as \textbf{< 5}, and any rate built on it as \textbf{—}, so that a
narrow filter cannot point at one person. A zero is shown, because it points at nobody. The
threshold is a setting of the deployment; ask whoever runs it if your site uses another. The page
does not hide a count that could be worked out from a total and its neighbors, so treat a
small department or week with care when you pass the figures on.
\end{elicitnote}

The participation, the weeks, each choice question and the progress by step are also drawn as
charts beside their tables. A count shown as \textbf{< 5} draws no bar.

\ui{Download CSV} and \ui{Download PDF} save the tables the page is showing (the PDF with the
charts too), with the filter at the
top and the same suppressed counts. They hold no respondent's record; that is what the export of
Chapter~\ref{sec:move} is for.

\begin{elicitnote}[Where the marks come from]
The marked questions are part of the survey definition: the author marks them in the authoring
tool and they arrive when the survey is applied (Chapter~\ref{sec:install}). A survey with no marks shows the
participation tables and says that nothing is marked.
\end{elicitnote}

\elicitpart{Surveys and the deployment}{Installing and updating surveys, moving respondents
between deployments, and the screens that report on the deployment itself, for the
administrator; then, for everyone, reading the console in your own language.}

%% ------------------------------------------------------------------ 11
\section{Moving a respondent between deployments}\adminonly
\label{sec:move}

A respondent can be moved from one Elicit deployment to another — a site sending its finished
respondents to the study's master site, say. The respondent is exported to a file at one end and
imported from it at the other. Both deployments must have the same survey installed.

\subsection{Exporting}

\begin{enumerate}
  \item Open \ui{Search Subjects} and find the respondent.
  \item In their row's \ui{Action} column, choose \ui{Export}, then \ui{Submit}.
\end{enumerate}

\screenshot{16-respondent-export}{The \ui{Export} action. It is offered to administrators
only, which is why a user does not see it on the same row.}

The browser downloads a \file{.elicit} file holding that one respondent: their answers, their
subject record, the code of the department they belong to, their access code, and whether they
have finished. It does not hold the survey definition.

\subsection{Importing}

\begin{enumerate}
  \item Open \ui{Admin} $>$ \ui{Import Respondent} on the receiving deployment.
  \item Choose \ui{Upload} and select the \file{.elicit} file.
\end{enumerate}

\screenshot{17-respondent-import}{The import screen. One respondent file at a time; the
result says how many records were imported, or why it was refused.}

Three things must be true on the receiving deployment, and the import is refused with a reason
when they are not:

\begin{itemize}
  \item the same survey must be installed — the file carries the survey's key and is matched
        against it;
  \item a department with the same \emph{code} must already exist. The import never creates one,
        so create it first (Chapter~\ref{sec:departments}) with the code the sending site used;
  \item the access code in the file must not already be in use here. It is never quietly
        reissued.
\end{itemize}

An imported respondent arrives as they left: a finished respondent is still finished, and their
answers read as they answered them.

%% ------------------------------------------------------------------ 12
\section{Installing and updating a survey}\adminonly
\label{sec:install}

A survey is built in the Author tool and arrives here as a \file{.elicit} definition file. The
same file installs it on a deployment that has never seen it and updates one that has — you do
not have to say which.

\subsection{Applying a definition}

\begin{enumerate}
  \item Open \ui{Admin} $>$ \ui{Apply Survey Definition}.
  \item Choose \ui{Upload} and select the \file{.elicit} file.
\end{enumerate}

\screenshot{18-apply-survey-definition}{The apply screen. The file identifies itself, so
this deployment works out whether it is a new survey or a revision of one already here.}

\screenshot{19-apply-result}{The result of an apply. \ui{New Survey Installed} means this
deployment had never seen this survey; a revision of one already here reports \ui{Survey
Updated}.}

The result names the survey key, the number of records installed and a breakdown by kind. A new
survey is in no department yet, and the result says so: open \ui{Departments} and add it to each
department that runs it before anyone registers onto it (Chapter~\ref{sec:departments}). When
it ends with \emph{Reporting schema rebuilt}, the reporting tables were regenerated to match the
new structure as part of the same apply — nothing needs restarting.

\begin{elicitcaution}
A revision older than the one already installed is refused. Going back to an earlier revision is
a database restore, not an upload.
\end{elicitcaution}

Two other endings need someone's attention:

\begin{description}
  \item[\ui{Apply Failed}] Nothing was installed. The message says why, most often a file that
        is not a survey definition or was exported by an older authoring tool.
  \item[\emph{Reporting schema not rebuilt}] The survey is installed, but its reporting tables
        were not brought up to date; the line gives the reason. Apply the same file again once
        the cause is fixed, or ask whoever runs the deployment to restart Survey. When the
        deployment builds reporting tables only at a restart, the result says nothing about
        them at all.
\end{description}

\subsection{Exporting a definition}

\ui{Admin} $>$ \ui{Export Survey Definition} downloads the definition of any survey installed
here. The file carries the survey's structure only — never respondent data — and can be applied
to another deployment or opened in the Author tool.

\screenshot{20-export-survey-definition}{Every installed survey, with its key, the
revision this deployment is running and the name of its reporting schema. \ui{Download} takes
the definition as it stands now; \ui{Rename} is described below.}

Each download is a fresh revision of the definition as it stands at that moment. To roll one
revision out across several sites, download it once and distribute \emph{that} file, rather than
exporting again at each site.

\subsection{Renaming a survey's reporting schema}

Every installed survey has a reporting schema of its own in the database, where its answers are
laid out for analysts and reporting tools. The schema is named after the survey when the survey
is first applied — \val{report_family_history_survey} for the Family Health History Survey —
and the \ui{Export Survey Definition} list shows the name beside each survey. A site that wants a
shorter name, or one that follows its own convention, renames it here. Nothing in the schema is
lost or rebuilt: the tables, views and every row keep going under the new name, and the console,
Survey and the report services pick the new name up at once.

\begin{enumerate}
  \item Open \ui{Admin} $>$ \ui{Export Survey Definition}.
  \item Choose \ui{Rename} on the survey's row.
  \item Type the new name: 1 to 63 lower-case letters, digits and underscores, not starting
        with a digit. The field says so while the name is not acceptable.
  \item Choose \ui{Rename} in the dialog. The list shows the new name.
\end{enumerate}

\screenshot{30-rename-reporting-schema}{The rename dialog: the survey, its current schema
name, the caution about what is outside the console, and the new name. \ui{Rename} stays
disabled until the name is acceptable.}

\begin{elicitcaution}
Queries, saved reports and BI connections \emph{outside} Elicit that name the current schema
stop working after the rename; whoever maintains them needs the new name. Inside Elicit nothing
is affected. A name another survey already uses is refused, and so is a name the database
reserves.
\end{elicitcaution}

A survey whose row says \emph{Not built yet} has no reporting schema: its definition was
installed but the reporting tables have not been built. Applying the definition again builds
them and names the schema, after which \ui{Rename} is offered.

\subsection{Before any survey is installed}

A deployment with no survey says so on every screen, and nobody can be registered until one is
applied.

\screenshot{21-missing-survey-notice}{The notice a deployment carries until a survey
definition is applied. An administrator is offered the remedy as a link; a user is told to ask
for one.}

%% ------------------------------------------------------------------ 13
\section{The System screens}\adminonly
\label{sec:system}

The six \ui{System} screens report what this deployment is running and test that the things it
depends on answer. None of them changes any configuration: the settings they show were fixed
when the service started, and they are shown as present or absent rather than by value.

\subsection{Overview}

\screenshot{22-system-overview}{\ui{System} $>$ \ui{Overview}: what is running, whether it is
healthy, and the setup work still outstanding.}

\ui{Running} gives the application, its version, when it was built and how long it has been up.
\ui{Health} gives each area's state with the detail behind it. \ui{Required settings} lists the
five settings the deployment cannot run without (the two database passwords, the OIDC client
secret, the mail sender address and the service token) and whether each was supplied.
\ui{Setup still outstanding} is the one to read first on a new installation, and links each item
to the screen that fixes it. It names:

\begin{itemize}
  \item no department, no message template, no template carrying the access code placeholder,
        no department with a sender address, and no survey;
  \item a service token still at its development value;
  \item a setting that is configured but no longer read, with the one that replaced it;
  \item a request to show identity tokens that the deployment ignores outside development.
\end{itemize}

\subsection{Database}

\screenshot{23-system-database}{\ui{System} $>$ \ui{Database}: both connections, the
migrations each module has applied, and what the database holds.}

The application connection serves the console; the owner connection runs the migrations. The
readiness probe only tests the first, so a wrong owner password shows up here and nowhere else.

\subsection{Branding}

\screenshot{24-system-branding}{\ui{System} $>$ \ui{Branding}: where each brand asset was
resolved from.}

A brand resolves from the mounted directory first, then the local directory, then the default
packaged with the application. Each asset is listed separately because a partial mount renders
with the wrong fonts and raises no error. \ui{Reload brand} discards the cache and resolves
again.

\subsection{Email}

\screenshot{25-system-email}{\ui{System} $>$ \ui{Email}: the mail settings the service
started with, and a test message.}

\ui{Send test email}, under \ui{Send a test message}, sends through the same service, relay and timeout that invitations use,
so a test that passes means the relay accepts mail. The test is sent from the service's own sender
address, not from a department's \ui{From email}, so it does not prove that the relay accepts each
department's address. Sending is disabled when no sender address is configured.

\subsection{Connections}

\screenshot{26-system-connections}{\ui{System} $>$ \ui{Connections}: every outbound
dependency, checked one at a time or all at once.}

Each check sends one read-only request and gives up after a few seconds. Post-survey actions are
never invoked and no report is generated, so checking is safe at any time.

\subsection{OIDC}

\screenshot{27-system-oidc}{\ui{System} $>$ \ui{OIDC}: the identity the console is currently
running under.}

This screen answers "why can this account not see that screen?". It shows the signed-in username,
the roles that reached the console, where those roles came from, and the identity provider's
tokens themselves (masked unless the deployment is in development). These are sign-in tokens,
nothing to do with a respondent's access code.

\subsection{The seeded-account warning}

\screenshot{28-default-account-warning}{The warning every screen carries until the seeded
accounts are renamed.}

The warning stays on every screen until the two seeded accounts are renamed
(Chapter~\ref{sec:users}); \ui{Open Users} takes you there.

%% ------------------------------------------------------------------ 14
\section{Reading the console in your language}
\label{sec:language}

The console starts in the language your browser prefers, when the deployment offers it, and
in English otherwise. When a deployment has more than one language available, a selector
appears at the right of the header. The choice lasts for your browser session only; nobody
else's console changes.

\begin{enumerate}
  \item Choose the selector in the header.
  \item Choose a language.
\end{enumerate}

\screenshot{29-language-selector}{The language selector, open. The tick marks the language
in use; the list holds every language this deployment has available.}

The console reloads in the chosen language, laid out right to left for a language such as
Arabic. A text the chosen language does not translate is shown in English rather than left
blank, so a partly translated console is still usable.

\begin{elicitnote}[Only English?]
The selector is hidden when English is the only language available, because there would be
nothing to choose. Which languages a deployment offers is set when it is installed, not from the
console: the console's image carries English, and every other language is a translation file the
installer mounts beside it, as the installation manual describes. Ask whoever runs your
deployment.
\end{elicitnote}

The survey's own questions are not translated by this selector. What a respondent reads is part
of the survey definition and is the author's to write.

%% ------------------------------------------------------------------ 15
\section{Where to look next}
\label{sec:next}

\begin{itemize}
  \item The \textbf{installation manual} covers everything this one does not: bringing up the
        database, the identity provider and the two services, the configuration reference,
        provisioning a brand, mounting the translations of the languages this deployment
        offers, and building and installing report services.
  \item The \textbf{author's manual}, \emph{Building a Survey}, covers the Author tool — writing
        questions, rules and reports, and exporting the definition that Chapter~\ref{sec:install} applies.
  \item \ui{System} $>$ \ui{Overview} is the first place to look when something is wrong here.
        It names what is unhealthy and what setup is still outstanding, and links to the screen
        that fixes it.
\end{itemize}

\end{document}
